\chapter*{KẾT LUẬN VÀ HƯỚNG PHÁT TRIỂN}
\addcontentsline{toc}{chapter}{KẾT LUẬN VÀ HƯỚNG PHÁT TRIỂN}

\section*{1. Kết quả đạt được}
\addcontentsline{toc}{section}{1. Kết quả đạt được}

\subsection*{1.1. Về thiết kế cơ sở dữ liệu}
\addcontentsline{toc}{subsection}{1.1. Về thiết kế cơ sở dữ liệu}

Đề tài đã phân tích bài toán thư viện số tư liệu nghe nhìn dùng chung cho nhiều đơn vị và
thiết kế một cơ sở dữ liệu quan hệ gồm 30 bảng chia thành sáu nhóm: danh mục thư viện, kho dữ
liệu video, chỉ mục văn bản, tài khoản và phân quyền, yêu cầu khai thác và cấp phép, nhật ký
và đánh giá. Lược đồ có đủ các loại liên kết: một--nhiều (đơn vị -- bộ sưu tập -- video),
nhiều--nhiều qua bảng trung gian (video -- chủ đề, vai trò -- quyền) và quan hệ đệ quy (cây
chủ đề). Lược đồ được chuẩn hoá có nêu từng bước: tách thuộc tính đa trị để đạt 1NF, kiểm tra
phụ thuộc bộ phận cho 2NF, loại phụ thuộc bắc cầu (đơn vị sở hữu suy ra qua bộ sưu tập, văn
bản bỏ dấu là cột tính toán) để đạt 3NF. Kết quả cuối cùng đạt BCNF; hai dư thừa giữa các bảng
được giữ có chủ đích, được bảo vệ bằng giao dịch và quyền, và có script kiểm tra nhất quán.

\subsection*{1.2. Về cài đặt trên SQL Server}
\addcontentsline{toc}{subsection}{1.2. Về cài đặt trên SQL Server}

Toàn bộ thiết kế được cài đặt bằng T-SQL và chạy trên SQL Server Express, gồm:
\begin{itemize}
    \item 128 ràng buộc có tên (30 PK, 37 FK, 13 UQ, 24 CK, 24 DF) và 109 cột NOT NULL;
    \item 5 hàm, trong đó có một hàm tất định dùng cho cột tính toán và một hàm bảng dùng CTE
    đệ quy;
    \item 10 thủ tục lưu trữ, có giao dịch tường minh, khoá \texttt{UPDLOCK} chống duyệt trùng,
    \texttt{SEQUENCE} sinh mã, tham số OUTPUT, và một thủ tục tự đăng ký khoá cứng quyền ngay
    trong CSDL;
    \item 11 view, có một view cập nhật được với \texttt{WITH CHECK OPTION};
    \item 6 trigger ghi nhật ký;
    \item phân quyền hai tầng: theo vai trò nghiệp vụ (bạn đọc, cán bộ thư viện, quản trị)
    trong ứng dụng và ngay trong thủ tục, và bằng role, người dùng không login cùng
    \texttt{EXECUTE AS} ở mức CSDL.
\end{itemize}
Vì bản Express không có Full-Text Search, đề tài tự xây chỉ mục ngược bằng bảng. Chỉ mục này
giảm số trang phải đọc từ 43 xuống 4 so với \texttt{LIKE} và chỉ khớp nguyên từ.

\subsection*{1.3. Về thực nghiệm}
\addcontentsline{toc}{subsection}{1.3. Về thực nghiệm}

Một tệp script dựng lại toàn bộ CSDL cùng dữ liệu mẫu (25 video và danh mục minh hoạ) trong
khoảng 5 giây. 25 truy vấn nghiệp vụ và bộ kiểm thử ràng buộc, giao dịch, phân quyền, chỉ mục,
nhất quán chạy không lỗi; 67 kiểm thử tự động đều qua. Trên 120 yêu cầu khai thác mẫu đã xác
minh, gộp ba nguồn (hình ảnh, chữ, lời nói) đạt Recall@100 = 0,525, cao hơn từng nguồn riêng
lẻ (0,358 và 0,308). Mọi lượt tìm xong dưới 0,6 giây, đáp ứng yêu cầu 30 giây. Ứng dụng
có giao diện lưu trữ 3D với ba cách trình bày kết quả, bộ lọc theo cây chủ đề và bộ sưu tập
dùng hàm đệ quy của CSDL, và trang tự đăng ký tài khoản bạn đọc.

\section*{2. Hạn chế}
\addcontentsline{toc}{section}{2. Hạn chế}

\begin{itemize}
    \item Đề tài dùng tình huống giả định. Video lấy từ một bộ dữ liệu công khai; các đơn vị
    thành viên khác, bạn đọc, yêu cầu và giấy phép là dữ liệu minh hoạ, chưa thử với đơn vị và
    người dùng thật.
    \item Chỉ một trong năm bộ sưu tập có video (lô L23, khoảng 1,5\% bộ dữ liệu), và 120 yêu
    cầu mẫu đều thuộc cùng một chủ đề (đua xe đạp), nên các chỉ số đánh giá còn dao động lớn.
    \item Mức truy cập của bộ sưu tập (mở/hạn chế) mới được lưu và thống kê, chưa được dùng để
    chặn việc xem hay xuất tư liệu khi chưa có giấy phép.
    \item Gộp kết quả với trọng số bằng nhau làm giảm Recall@1 so với chỉ dùng hình ảnh.
    \item Kênh hình ảnh nằm ngoài CSDL vì SQL Server 2017--2022 không có kiểu dữ liệu vector.
    \item Chỉ mục ngược và nhãn chủ đề được lập lại theo lô chứ chưa cập nhật tức thời khi kho
    nhận video mới.
    \item Cờ bật/tắt tự đăng ký chỉ được kiểm tra ở tầng web; thủ tục
    \texttt{sp\_register\_reader} vẫn gọi được khi cờ tắt, dù chỉ tạo được tài khoản bạn đọc.
    \item Giao diện 3D tải ảnh khung hình gốc (khoảng 170 KB mỗi ảnh) cho cả các ô nhỏ, nên chế
    độ Constellation cần 10--20 giây để hiện đủ 50 kết quả.
\end{itemize}

\section*{3. Hướng phát triển}
\addcontentsline{toc}{section}{3. Hướng phát triển}

\begin{itemize}
    \item Dùng mức truy cập của bộ sưu tập trong thủ tục: bạn đọc chỉ xuất được đoạn tư liệu
    thuộc bộ sưu tập hạn chế khi có giấy phép còn hiệu lực.
    \item Sinh sẵn ảnh thu nhỏ cho từng khung hình để giao diện 3D và trang chi tiết tải nhanh
    hơn.
    \item Thêm phạm vi giấy phép theo từng đoạn tư liệu, gia hạn và thu hồi giấy phép, kèm
    thông báo cho bạn đọc trước ngày hết hạn.
    \item Lưu luật gắn chủ đề trong bảng để cán bộ thư viện chỉnh sửa, và để chủ đề được gắn
    tự động khi nạp video mới.
    \item Lưu trọng số gộp của từng nguồn trong bảng \texttt{sources} để điều chỉnh không cần
    sửa mã, và chọn trọng số trên một tập yêu cầu mẫu riêng.
    \item Mở rộng sang toàn bộ bộ dữ liệu (401.153 khung hình) với phân vùng bảng theo lô, và
    dùng kiểu dữ liệu vector của các phiên bản SQL Server mới để đưa cả kênh hình ảnh vào CSDL.
\end{itemize}
